\chapter{Glosario}

\newcommand{\gl}[2]{\par\smallskip\textbf{\color{azul}#1.}\ #2}

\gl{Apalancamiento}{Exposición a ganancias y pérdidas mucho mayores que el capital comprometido. Cap.~4.}
\gl{Arbitraje}{Ganancia sin riesgo ni inversión neta por precios inconsistentes de activos equivalentes. Cap.~4.}
\gl{Backwardation}{Estructura en la que los futuros más lejanos son más baratos que los cercanos; refleja escasez y un rendimiento de conveniencia alto. Cap.~6.}
\gl{Base}{Diferencia entre el precio del activo cubierto y el del futuro usado como cobertura. Riesgo de base: incertidumbre sobre esa diferencia. Cap.~6.}
\gl{Basis (entre monedas)}{Diferencia entre el costo de una moneda en el mercado local de forwards y el implícito en las curvas internacionales. Cap.~7.}
\gl{Bootstrap}{Procedimiento que obtiene factores de descuento sucesivos a partir de tasas par (swap). Cap.~2.}
\gl{Call}{Opción que da el derecho de comprar a un precio fijado. Cap.~8.}
\gl{Capa}{Tramo de un programa rodante con su propio vencimiento; en el informe, nov, dic y ene, de 0.774 MMbl cada una. Cap.~12.}
\gl{Capital de trabajo}{Activo corriente menos pasivo corriente. Cap.~1.}
\gl{CCS (cross currency swap)}{Swap que intercambia intereses y principal en dos monedas. Cap.~7.}
\gl{Colateral}{Garantía (normalmente efectivo) que entrega la parte con valor negativo en un derivado, según el CSA. Cap.~4 y 13.}
\gl{Collar}{Put comprada más call vendida con strike mayor; de costo cero si las primas se igualan. Cap.~8.}
\gl{Contango}{Estructura en la que los futuros más lejanos son más caros que los cercanos. Cap.~6.}
\gl{Costo de acarreo}{Costo neto de mantener un activo hasta una fecha: intereses y almacenamiento menos rendimiento de conveniencia; $F=Se^{(r+u-y)T}$. Cap.~6.}
\gl{Covenant}{Obligación contractual de un deudor cuyo incumplimiento permite exigir el pago anticipado. Cap.~10.}
\gl{Crack spread}{Margen de refinación: precio de los productos menos el del crudo. Cap.~10.}
\gl{Cross-acceleration}{Versión suave del incumplimiento cruzado: solo se activa si otra deuda fue efectivamente acelerada. Cap.~4.}
\gl{Cross-default}{Cláusula que permite terminar los derivados si la contraparte incumple otra deuda. Cap.~4.}
\gl{CSA}{Anexo del contrato ISDA que regula el colateral. Cap.~4.}
\gl{CVA}{Ajuste del valor de un derivado por el riesgo de incumplimiento de la contraparte. Cap.~4.}
\gl{Delta}{Sensibilidad del valor de una opción al precio del subyacente. Cap.~9.}
\gl{Derivado}{Contrato cuyo valor depende del de otra variable (subyacente). Cap.~1.}
\gl{Derivado hipotético}{Derivado imaginario que replica exactamente la partida cubierta; sirve para medir la inefectividad. Cap.~13.}
\gl{Duración modificada}{Cambio porcentual aproximado del precio de un bono por cada unidad de cambio del rendimiento. Cap.~3.}
\gl{DVA}{Ajuste por el riesgo de crédito propio en el valor razonable de un derivado (NIIF~13). Cap.~4.}
\gl{EPE}{Exposición esperada positiva: promedio de lo que una parte perdería si la otra incumpliera. Cap.~4.}
\gl{Factor de descuento}{Valor hoy de una unidad monetaria pagada en una fecha futura. Cap.~2.}
\gl{FEPC}{Fondo de Estabilización de Precios de los Combustibles. Cap.~1 y 10.}
\gl{Forward}{Acuerdo OTC de compraventa futura a precio fijado hoy. Cap.~5.}
\gl{Futuro}{Forward estandarizado negociado en bolsa con liquidación diaria. Cap.~6.}
\gl{FVA y KVA}{Ajustes de valoración por el costo de fondeo y de capital del banco. Cap.~4.}
\gl{Gamma}{Sensibilidad de la delta al precio del subyacente. Cap.~9.}
\gl{ISDA}{Contrato marco estándar para derivados OTC. Cap.~4.}
\gl{Margen}{Depósito de garantía exigido por una bolsa; las llamadas de margen reponen pérdidas diarias. Cap.~4.}
\gl{Moneda funcional}{Moneda del entorno económico principal de la empresa; para PETROPERÚ, el dólar. Cap.~1.}
\gl{MTA}{Monto mínimo de transferencia de colateral en un CSA. Cap.~4.}
\gl{NDF}{Forward de monedas que se liquida por diferencias, sin entrega del principal. Cap.~5.}
\gl{NIIF 9}{Norma internacional sobre instrumentos financieros; su capítulo 6 regula la contabilidad de coberturas. Cap.~13.}
\gl{No arbitraje}{Principio según el cual activos con iguales flujos futuros deben tener igual precio. Cap.~4 y 5.}
\gl{Nocional}{Monto de referencia sobre el que se calculan los flujos de un derivado. Cap.~7.}
\gl{ORI}{Otro resultado integral: cuenta del patrimonio donde se difiere la parte efectiva de una cobertura de flujos. Cap.~13.}
\gl{OTC}{Mercado extrabursátil, de contratos negociados directamente entre las partes. Cap.~4.}
\gl{Paridad cubierta de tasas}{$F=S_0[(1+r_{dom})/(1+r_{ext})]^T$: el forward de monedas refleja el diferencial de tasas. Cap.~5.}
\gl{Paridad put-call}{Relación de no arbitraje entre call y put con igual strike y vencimiento. Cap.~8.}
\gl{Payoff}{Pago contractual de un derivado al vencimiento. Cap.~4.}
\gl{PCC, PND, PCG}{Posición de cambio contable, en derivados y global ($\text{PCG}=\text{PCC}+\text{PND}$). Cap.~11.}
\gl{PFE}{Exposición potencial futura: cuantil alto (95\,\% o 99\,\%) de la exposición en una fecha futura. Cap.~4 y 13.}
\gl{Pip}{Una diezmilésima de unidad en un tipo de cambio (0.0001). Cap.~5.}
\gl{Posición larga / corta}{Gana si el precio sube / si baja. Cap.~4.}
\gl{Prima}{Precio de una opción. Cap.~8.}
\gl{Programa rodante}{Cobertura de un riesgo permanente con contratos de corto plazo escalonados que se reponen al vencer. Cap.~12.}
\gl{Put}{Opción que da el derecho de vender a un precio fijado. Cap.~8.}
\gl{Rate lock}{Contrato que fija hoy la tasa base de una emisión o préstamo futuro. Cap.~7.}
\gl{Rendimiento de conveniencia}{Beneficio implícito de tener físicamente una materia prima. Cap.~6.}
\gl{Saldo estructural}{Parte del inventario que siempre está presente aunque sus unidades roten. Cap.~12.}
\gl{Skew (sesgo de volatilidad)}{Diferencia de volatilidad implícita entre opciones con distinto strike. Cap.~9.}
\gl{SOFR}{Tasa de referencia de préstamos a un día garantizados con bonos del Tesoro de EE.\,UU. Cap.~2.}
\gl{Spread de crédito}{Rendimiento adicional exigido a un emisor sobre la tasa sin riesgo. Cap.~3.}
\gl{Strike}{Precio de ejercicio de una opción. Cap.~8.}
\gl{Swap}{Intercambio de series de flujos calculados con reglas distintas. Cap.~7.}
\gl{Swap de inicio diferido}{Swap cuyos flujos empiezan en una fecha futura con tasa fijada hoy. Cap.~7.}
\gl{Tasa cupón cero}{Rendimiento de una inversión con un único pago al vencimiento. Cap.~2.}
\gl{Theta}{Cambio del valor de una opción por el paso del tiempo. Cap.~9.}
\gl{Umbral (threshold)}{Pérdida que un banco tolera en un derivado antes de pedir colateral. Cap.~4 y 13.}
\gl{Valor razonable}{Valor de mercado de un activo o pasivo. Cap.~5 y 13.}
\gl{VaR}{Valor en riesgo: pérdida que solo se supera con una probabilidad dada (por ejemplo, 5\,\%) en un horizonte. Cap.~3.}
\gl{Vega}{Sensibilidad del valor de una opción a la volatilidad. Cap.~9.}
\gl{Volatilidad}{Desviación estándar anualizada de los rendimientos logarítmicos; histórica o implícita. Cap.~3.}
\gl{Wrong-way}{Riesgo de que la exposición frente a una contraparte crezca justo cuando ella se debilita. Cap.~4.}
\gl{WTI}{West Texas Intermediate: crudo de referencia negociado en NYMEX. Cap.~1 y 6.}
